\documentclass[toc=bibliographynumbered]{scrartcl} %mit [draft] vor den {} erreicht man, daß statt der Grafiken nur ein Rahmen in Größe der Grafik gezeichnet wird. Zusätzlich wird der Name der einzufügenden Datei in diesen Rahmen geschrieben]{book}
\usepackage[utf8]{inputenc}
\usepackage[T1]{fontenc}
\usepackage{lmodern}
%\usepackage[english,ngerman]{babel} %Deutsch als Hauptsprache, z.b. für Inhaltsverzeichnis etc.
\usepackage[ngerman,english]{babel} % Englisch als Hauptsprache!!!

\usepackage{graphicx}
\usepackage{xspace}
\usepackage{amstext}
\usepackage{amssymb}


\usepackage[paper=a4paper,margin=2.5cm]{geometry}
\usepackage{helvet}
\renewcommand*{\familydefault}{\sfdefault} %gehört zum Font


\usepackage{here} %erzwingt Platzierung von Bildern ("H" bei Option)


\begin{document}

\begin{titlepage}
%\begin{figure}[t]
\begin{center}
\includegraphics[scale=0.65]{pics/hulogo.pdf}\\[1cm]

%\end{figure}
\textsc{\textbf{\normalsize LEBENSWISSENSCHAFTLICHE FAKULTÄT\\[0.2cm]  INSTITUT FÜR BIOLOGIE\\[1.5cm]
\LARGE MASTERARBEIT\\[0.5cm]
\normalsize ZUM ERWERB DES AKADEMISCHEN GRADES\\[0.5cm] MASTER OF SCIENCE\\[2cm]}
\normalfont %Kapitälchen deaktivieren
\textit{\large "Körpergrößentrends in fossilen Landschildkröten aus dem Neogen"\\[0.5cm]
"Body size trends in Neogene testudinid tortoises" }\\[2cm]
\small vorgelegt von\\[0.5cm] Julia Joos \\[0.2cm] geb. am 18.05.1991 in Freudenstadt \\[1cm] angefertigt in der Arbeitsgruppe Paläozoologie\\[0.2cm]
am Institut für Biologie/Museum für Naturkunde\\[1cm]
Berlin, im Septenber 2017}
\end{center}


\end{titlepage}




\end{document}